\documentclass[a4paper,12pt]{article}
\usepackage[utf8]{inputenc}
\usepackage[T1]{fontenc}
\usepackage[ngerman]{babel}
\usepackage{amsmath,amssymb}
\usepackage{booktabs}
\usepackage{enumitem}
\usepackage[margin=2.5cm]{geometry}
\usepackage{tcolorbox}

\tcbuselibrary{breakable}

\newtcolorbox{merkbox}[1][]{
  colback=blue!5,
  colframe=blue!40,
  fonttitle=\bfseries,
  title={Merke},
  breakable,
  #1
}

\newtcolorbox{analogie}[1][]{
  colback=green!5,
  colframe=green!40,
  fonttitle=\bfseries,
  title={Analogie},
  breakable,
  #1
}

\begin{document}

\begin{center}
{\LARGE\bfseries \"Ubungsblatt 3 -- Erkl\"arungen}\\[0.5em]
{\large VC Logik}
\end{center}

\bigskip

% ============================================================
% TEIL 0: Grundlagen
% ============================================================

\section*{Teil 0: Grundlagen -- alles, was man wissen muss}

Bevor wir uns an die Aufgaben setzen, lernen wir in aller Ruhe die
Begriffe, die im ganzen \"Ubungsblatt vorkommen. Keine Panik: jedes
Wort, jedes Symbol wird erkl\"art. Wir tun so, als h\"atten wir noch
nie etwas von Logik geh\"ort. Am Ende von Teil~0 bist du bestens
ger\"ustet f\"ur die drei Aufgaben.

\subsection*{0.1 Die f\"unf Symbole der Aussagenlogik}

In der Logik verwenden wir nicht lange S\"atze, sondern kleine Symbole,
damit wir schnell und eindeutig schreiben k\"onnen. Es gibt f\"unf
Symbole, die in diesem \"Ubungsblatt immer wieder vorkommen. Wir
schauen sie uns einzeln an.

\begin{center}
\renewcommand{\arraystretch}{1.5}
\begin{tabular}{c|l|l}
\toprule
\textbf{Symbol} & \textbf{Name} & \textbf{Bedeutung in Alltagssprache} \\
\midrule
$\neg A$ & Negation (Nicht) & "`$A$ ist falsch"' / "`Nicht $A$"' \\
$A \wedge B$ & Konjunktion (Und) & "`Beide, $A$ und $B$, sind wahr"' \\
$A \vee B$ & Disjunktion (Oder) & "`Mindestens eines von beiden ist wahr"' \\
$A \to B$ & Implikation (Wenn-Dann) & "`Wenn $A$ wahr ist, dann auch $B$"' \\
$A \leftrightarrow B$ & Biimplikation (Genau dann) & "`$A$ und $B$ haben denselben Wert"' \\
\bottomrule
\end{tabular}
\end{center}

\begin{merkbox}
In der Logik gibt es nur zwei Werte:
\begin{itemize}[nosep]
\item $1$ = \textbf{wahr} (die Aussage stimmt)
\item $0$ = \textbf{falsch} (die Aussage stimmt nicht)
\end{itemize}
Jede Variable, jede Formel hat am Ende einen dieser beiden Werte.
Es gibt kein "`vielleicht"'.
\end{merkbox}

\subsubsection*{Wahrheitstabelle f\"ur jedes Symbol einzeln}

Jedes Symbol hat eine kleine Tabelle, die festlegt, wie es rechnet.
Diese Tabellen muss man im Kopf haben -- ohne sie kommt man nirgendwo
hin.

\textbf{Negation $\neg A$} (dreht den Wert um):

\begin{center}
\renewcommand{\arraystretch}{1.2}
\begin{tabular}{c|c}
\toprule
$A$ & $\neg A$ \\
\midrule
0 & 1 \\
1 & 0 \\
\bottomrule
\end{tabular}
\end{center}

Analogie: Ein Lichtschalter mit Umkehrfunktion. Ist das Licht an,
macht die Negation es aus. Ist es aus, macht sie es an.

\textbf{Konjunktion $A \wedge B$} (wahr nur wenn beide wahr sind):

\begin{center}
\renewcommand{\arraystretch}{1.2}
\begin{tabular}{cc|c}
\toprule
$A$ & $B$ & $A \wedge B$ \\
\midrule
0 & 0 & 0 \\
0 & 1 & 0 \\
1 & 0 & 0 \\
1 & 1 & 1 \\
\bottomrule
\end{tabular}
\end{center}

\begin{analogie}
"`Ich darf fernsehen \textbf{und} ich habe meine Hausaufgaben
gemacht."' Nur wenn \emph{beide} Aussagen zutreffen, stimmt der ganze
Satz. Fehlt einer der beiden Teile, ist das Ganze falsch -- dann sagt
Mama "`kein Fernsehen"'.
\end{analogie}

\textbf{Disjunktion $A \vee B$} (wahr wenn mindestens eines wahr ist):

\begin{center}
\renewcommand{\arraystretch}{1.2}
\begin{tabular}{cc|c}
\toprule
$A$ & $B$ & $A \vee B$ \\
\midrule
0 & 0 & 0 \\
0 & 1 & 1 \\
1 & 0 & 1 \\
1 & 1 & 1 \\
\bottomrule
\end{tabular}
\end{center}

\begin{analogie}
"`Zum Mittagessen gibt es Nudeln \textbf{oder} Pizza."' Wenn es Nudeln
gibt, stimmt der Satz. Wenn es Pizza gibt, stimmt er auch. Wenn es
\emph{beides} gibt (gro{\ss}er Hunger!), stimmt er ebenfalls. Nur wenn
es weder Nudeln noch Pizza gibt, ist die Aussage falsch. In der Logik
ist das "`Oder"' also ein \emph{einschlie{\ss}endes} Oder, nicht das
"`entweder...oder"' der Alltagssprache.
\end{analogie}

\textbf{Implikation $A \to B$} (nur falsch, wenn $A$ wahr und $B$ falsch):

\begin{center}
\renewcommand{\arraystretch}{1.2}
\begin{tabular}{cc|c}
\toprule
$A$ & $B$ & $A \to B$ \\
\midrule
0 & 0 & 1 \\
0 & 1 & 1 \\
1 & 0 & \textbf{0} \\
1 & 1 & 1 \\
\bottomrule
\end{tabular}
\end{center}

\begin{analogie}[title={Analogie: Regenversprechen}]
Dein Freund sagt: \textbf{"`Wenn es regnet, bringe ich einen Schirm
mit."'} Wann hat er gelogen?
\begin{itemize}[nosep]
\item Es regnet nicht, kein Schirm dabei $\Rightarrow$ \textbf{nicht
      gelogen}. Er hat ja nichts versprochen f\"ur "`kein Regen"'.
\item Es regnet nicht, aber er hat trotzdem einen Schirm
      $\Rightarrow$ \textbf{nicht gelogen}. Das ist h\"ochstens
      vorsichtig, aber kein Wortbruch.
\item Es regnet, aber er hat keinen Schirm $\Rightarrow$
      \textbf{gelogen!} Das war genau der Fall, den er versprochen
      hat abzudecken.
\item Es regnet, er hat einen Schirm $\Rightarrow$
      \textbf{nicht gelogen}. Versprechen gehalten.
\end{itemize}
Die Implikation ist also genau dann falsch, wenn die Voraussetzung
($A$, "`es regnet"') wahr ist und die Folge ($B$, "`Schirm dabei"')
falsch.
\end{analogie}

\begin{merkbox}
Die Implikation $A \to B$ ist \textbf{nur in einem einzigen Fall}
falsch: $A = 1$ und $B = 0$. In allen anderen drei F\"allen ist sie
wahr. Ein Spezialfall: Wenn $A = 0$, dann ist $A \to B$ automatisch
wahr -- egal, was $B$ ist. Das nennt man "`ex falso quodlibet"' oder
"`aus Falschem folgt alles"'.
\end{merkbox}

\textbf{Biimplikation $A \leftrightarrow B$} (wahr wenn beide
denselben Wert haben):

\begin{center}
\renewcommand{\arraystretch}{1.2}
\begin{tabular}{cc|c}
\toprule
$A$ & $B$ & $A \leftrightarrow B$ \\
\midrule
0 & 0 & 1 \\
0 & 1 & 0 \\
1 & 0 & 0 \\
1 & 1 & 1 \\
\bottomrule
\end{tabular}
\end{center}

\begin{analogie}
"`Ich gehe genau dann ins Schwimmbad, wenn die Sonne scheint."' Das
ist eine Biimplikation, weil sie \emph{beide} Richtungen behauptet:
(1)~Sonne $\to$ Schwimmbad, und (2)~Schwimmbad $\to$ Sonne. Du kannst
dich also nie im Schwimmbad befinden, ohne dass die Sonne scheint.
Biimplikation ist wahr, wenn beide Aussagen dasselbe tun -- entweder
beide wahr oder beide falsch.
\end{analogie}

\subsection*{0.2 Operatorpr\"azedenz -- wer bindet st\"arker?}

Jetzt kommt ein heikles Thema, das in Aufgabe~1 ganz wichtig wird:
Wenn in einer Formel keine Klammern stehen, in welcher Reihenfolge
werden dann die Operatoren ausgewertet? Zum Beispiel:
\[
\neg X \leftrightarrow Z \to \neg Y \wedge Z
\]
Was bindet hier zuerst? Ohne Regel k\"onnte man die Formel auf viele
verschiedene Arten lesen und w\"urde jedes Mal ein anderes Ergebnis
bekommen. Damit alle Menschen (und Computer) dieselbe Lesart haben,
gibt es feste Pr\"azedenzregeln:

\begin{merkbox}
\textbf{Operatorpr\"azedenz in der Aussagenlogik} (von stark nach schwach):
\begin{enumerate}[nosep]
\item $\neg$ (Negation) -- \textbf{st\"arkste} Bindung
\item $\wedge$ (Und)
\item $\vee$ (Oder)
\item $\to$ (Implikation)
\item $\leftrightarrow$ (Biimplikation) -- \textbf{schw\"achste} Bindung
\end{enumerate}
Eselsbr\"ucke: \textbf{N, U, O, W, G} -- \textbf{N}icht, \textbf{U}nd,
\textbf{O}der, \textbf{W}enn-dann, \textbf{G}enau-dann-wenn.
\end{merkbox}

\begin{analogie}[title={Analogie: Punkt-vor-Strich}]
In der Schule hast du gelernt: $3 + 4 \cdot 2 = 3 + 8 = 11$ und nicht
$14$, weil $\cdot$ st\"arker bindet als $+$. Die Mal-Rechnung wird
zuerst gemacht, und erst danach die Plus-Rechnung. In der Logik ist
es genauso: $\neg$ und $\wedge$ binden fest wie Mal-Rechnungen, $\to$
und $\leftrightarrow$ sind wie das lose $+$. Im Zweifel: Klammere
stillschweigend die "`stark"' gebundenen Operatoren zusammen, bevor du
die "`schwachen"' dranh\"angst.
\end{analogie}

\textbf{Beispiel.} Schauen wir uns nochmal die Formel aus Aufgabe~1 an:
\[
\neg X \leftrightarrow Z \to \neg Y \wedge Z.
\]
Nach der Pr\"azedenz klammern wir sie Schritt f\"ur Schritt:
\begin{enumerate}[nosep]
\item Am st\"arksten bindet $\neg$: aus $\neg X$ und $\neg Y$ werden
      Bausteine, die nichts mehr verlieren k\"onnen.
\item Dann $\wedge$: das Paar $\neg Y \wedge Z$ wird zu einem
      einzigen Baustein zusammengefasst.
\item Dann $\to$: jetzt bindet sich $Z$ mit dem schon zusammengesetzten
      $(\neg Y \wedge Z)$ zu $Z \to (\neg Y \wedge Z)$.
\item Zuletzt $\leftrightarrow$: das ganz oben stehende
      $\neg X$ wird mit dem rechten Baustein durch die Biimplikation
      verbunden.
\end{enumerate}
Das Endergebnis der Klammerung lautet:
\[
\neg X \;\leftrightarrow\; \bigl(\,Z \to (\neg Y \wedge Z)\,\bigr).
\]

\subsection*{0.3 Literal -- der kleinste Baustein}

\begin{merkbox}
Ein \textbf{Literal} ist entweder eine Variable oder ihre Negation.
Also: $A$ ist ein Literal. $\neg A$ ist auch ein Literal. Mehr nicht.
\end{merkbox}

\begin{analogie}
Stell dir ein Literal wie einen Legostein mit einem einzigen Knopf
vor. Entweder steht er "`an"' (die Variable) oder "`aus"' (die
negierte Variable). Gr\"o{\ss}ere Bauwerke (Formeln) entstehen aus
vielen solcher kleinen Bausteine.
\end{analogie}

\subsection*{0.4 Klausel und Minterm}

Jetzt bauen wir etwas gr\"o{\ss}ere Bauteile. Zwei Arten von
Zusammenbauten sind besonders wichtig:

\begin{merkbox}
\textbf{Klausel} = Disjunktion (ODER-Verkn\"upfung) von Literalen.\\
Beispiel: $(A \vee \neg B \vee C)$.

\medskip
\textbf{Minterm} = Konjunktion (UND-Verkn\"upfung) von Literalen.\\
Beispiel: $(A \wedge \neg B \wedge C)$.
\end{merkbox}

Der Unterschied ist, welcher Junktor im Inneren verwendet wird. Eine
Klausel verwendet nur $\vee$, ein Minterm nur $\wedge$.

\begin{analogie}
Eine \textbf{Klausel} ist wie eine Liste mit \textbf{Notausg\"angen}:
Mindestens einer muss offen sein, damit man rauskommt (mindestens ein
Literal muss wahr sein, damit die Klausel wahr ist).

\medskip
Ein \textbf{Minterm} ist wie eine Liste mit \textbf{T\"urcodes}: Alle
m\"ussen stimmen, sonst geht der Safe nicht auf (alle Literale
m\"ussen wahr sein, damit der Minterm wahr ist).
\end{analogie}

\subsection*{0.5 CNF -- Konjunktive Normalform}

\begin{merkbox}
\textbf{CNF (Konjunktive Normalform)} ist ein "`UND von ODERs"'. Die
Formel sieht aus wie:
\[
(L_1 \vee L_2 \vee \dots) \;\wedge\; (L_3 \vee L_4 \vee \dots)
  \;\wedge\; \dots
\]
Im Klartext: Eine Konjunktion (UND) von Klauseln. Jede Klausel ist
selbst eine Disjunktion (ODER) von Literalen.
\end{merkbox}

\begin{analogie}
CNF ist wie eine \textbf{Checkliste mit Anforderungen}. Jede Zeile
der Checkliste ist eine Klausel. Die ganze Formel ist nur dann wahr,
wenn \emph{jede} Anforderung erf\"ullt ist (UND). Eine einzelne
Anforderung ist erf\"ullt, wenn \emph{mindestens eine} ihrer
Alternativen zutrifft (ODER).

\medskip
Beispiel: "`Das Auto f\"ahrt, wenn (Diesel \textbf{oder} Benzin im
Tank) \textbf{und} (Batterie voll \textbf{oder} an der Steckdose)
\textbf{und} der Fahrer eingestiegen ist."' Alle drei Checks m\"ussen
bestanden werden, und jeder Check hat evtl.\ Alternativen.
\end{analogie}

\subsection*{0.6 DNF -- Disjunktive Normalform}

\begin{merkbox}
\textbf{DNF (Disjunktive Normalform)} ist ein "`ODER von UNDs"'. Die
Formel sieht aus wie:
\[
(L_1 \wedge L_2 \wedge \dots) \;\vee\; (L_3 \wedge L_4 \wedge \dots)
  \;\vee\; \dots
\]
Im Klartext: Eine Disjunktion (ODER) von Mintermen. Jeder Minterm ist
selbst eine Konjunktion (UND) von Literalen.
\end{merkbox}

\begin{analogie}
DNF ist wie eine \textbf{Liste von Erfolgsszenarien}. Jedes Szenario
(Minterm) beschreibt einen Zustand, in dem es klappt. Die ganze Formel
ist wahr, wenn mindestens \emph{ein} Szenario voll zutrifft (ODER),
und ein Szenario trifft zu, wenn \emph{alle} seine Bedingungen erf\"ullt
sind (UND).

\medskip
Beispiel: "`Ich gehe Essen, wenn (Mama kocht \textbf{und} es gibt
Nudeln) \textbf{oder} (Freund ruft an \textbf{und} bietet Pizza an)
\textbf{oder} (ich selbst habe Zeit \textbf{und} Geld)."' Jedes dieser
Szenarien reicht aus.
\end{analogie}

\begin{merkbox}
\textbf{Eselsbr\"ucke CNF vs.\ DNF.}
\begin{itemize}[nosep]
\item \textbf{D}NF = "`\textbf{D}as \textbf{N}ormale \textbf{F}ormat"'
      = Summe aus Produkten. (In der booleschen Algebra wird $\vee$
      oft als $+$ und $\wedge$ oft als $\cdot$ geschrieben, daher
      "`Summe von Produkten"'.)
\item \textbf{C}NF = wie sa\textbf{C}-gepackt = Produkt aus Summen.
      Die \"au{\ss}ere Struktur ist Multiplikation ($\wedge$), innen
      stehen Summen ($\vee$).
\end{itemize}
\end{merkbox}

\subsection*{0.7 Wahrheitstabelle -- die Landkarte einer Formel}

Eine Wahrheitstabelle listet f\"ur jede m\"ogliche Belegung der
Variablen den Wert der Formel auf. Bei $n$ Variablen gibt es $2^n$
Zeilen, weil jede Variable 2 M\"oglichkeiten hat und sich diese
multiplizieren.

\begin{center}
\renewcommand{\arraystretch}{1.3}
\begin{tabular}{c|c|l}
\toprule
\textbf{Variablen} & \textbf{Zeilen} & \textbf{Erkl\"arung} \\
\midrule
1 & 2 & $A = 0$ oder $A = 1$ \\
2 & 4 & Jede Kombination von $A$ und $B$ \\
3 & 8 & $2 \cdot 2 \cdot 2$ \\
4 & 16 & $2^4$ \\
$n$ & $2^n$ & Allgemein \\
\bottomrule
\end{tabular}
\end{center}

In Aufgabe~1 ist $n = 3$ (Variablen $X, Y, Z$), also brauchen wir
$8$~Zeilen. In Aufgabe~3 ist $n = 4$ (Variablen $A, B, C, D$), also
$16$~Zeilen. Die \"ublichste Reihenfolge ist das bin\"are Hochz\"ahlen
von $0\dots 0$ bis $1\dots 1$.

\subsection*{0.8 \"Aquivalenz ($\equiv$)}

\begin{merkbox}
Zwei Formeln $F$ und $G$ sind \textbf{\"aquivalent} ($F \equiv G$),
wenn sie \textbf{dieselbe Wahrheitstabelle} haben -- also f\"ur jede
Belegung denselben Wert liefern.
\end{merkbox}

Wichtig: $\equiv$ ist \emph{nicht} dasselbe wie $\leftrightarrow$!
\begin{itemize}[nosep]
\item $\leftrightarrow$ ist ein Operator \emph{innerhalb} einer Formel.
      Zum Beispiel: $A \leftrightarrow B$ ist selbst wieder eine
      Formel, die je nach Belegung wahr oder falsch sein kann.
\item $\equiv$ ist eine Aussage \emph{\"uber} zwei Formeln auf der
      Metaebene: "`Diese beiden Formeln berechnen dieselbe boolesche
      Funktion."'
\end{itemize}

\begin{analogie}
Stell dir zwei Kuchenrezepte vor. Wenn sie am Ende jedes Mal exakt
denselben Kuchen liefern, egal welche Zutaten verf\"ugbar sind, dann
sind die Rezepte \emph{\"aquivalent} -- auch wenn sie im Wortlaut
ganz unterschiedlich sind. Nicht "`dasselbe Rezept"', sondern
"`dasselbe Ergebnis"'. Genauso ist es mit Formeln und $\equiv$.
\end{analogie}

\subsection*{0.9 Wichtige \"Aquivalenzen zum Ausrechnen}

Wir brauchen ein paar Umformungsregeln, mit denen man Formeln
elegant umbauen kann. Diese werden sp\"ater, besonders in Aufgabe~2,
sehr n\"utzlich.

\begin{merkbox}
\textbf{Wichtige \"Aquivalenzen.}
\begin{align*}
\neg\neg A &\equiv A \quad &\text{(Doppelnegation)}\\
A \wedge B &\equiv \neg(\neg A \vee \neg B) \quad &\text{(De Morgan 1)}\\
A \vee B &\equiv \neg(\neg A \wedge \neg B) \quad &\text{(De Morgan 2)}\\
A \to B &\equiv \neg A \vee B \quad &\text{(Implikation als Oder)}\\
A \leftrightarrow B &\equiv (A \to B) \wedge (B \to A) \quad &\text{(Biimplikation zerlegen)}
\end{align*}
\end{merkbox}

\begin{analogie}[title={Analogie zu De Morgan}]
"`Nicht beide sind zu Hause"' ist dasselbe wie "`Mindestens einer ist
weg"'. Das ist De Morgan: Wenn ein UND verneint wird, wird es zu einem
ODER mit verneinten Einzelteilen. Und umgekehrt: "`Keiner ist zu
Hause"' ist dasselbe wie "`Sowohl $A$ ist weg als auch $B$ ist weg"'
-- aus einem verneinten ODER wird ein UND mit verneinten Teilen.
\end{analogie}

\subsection*{0.10 Aus Wahrheitstabelle zu CNF und DNF}

Jetzt kommt der \textbf{wichtigste Trick} f\"ur Aufgaben~1 und~3.
Wenn man die Wahrheitstabelle einer Formel hat, kann man daraus
mechanisch sowohl die CNF als auch die DNF ablesen.

\begin{merkbox}[title={Strategie: Wahrheitstabelle $\to$ DNF}]
\textbf{DNF-Rezept:}
\begin{enumerate}[nosep]
\item Gehe jede Zeile der Wahrheitstabelle durch, bei der die Formel
      den Wert $1$ hat ("`Einser-Zeilen"').
\item F\"ur jede solche Zeile bildest du einen Minterm (eine
      UND-Verkn\"upfung):
      \begin{itemize}[nosep]
      \item Ist eine Variable in der Zeile~$1$, \"ubernimm sie
            \emph{positiv} (z.\,B.\ $X$).
      \item Ist eine Variable in der Zeile~$0$, \"ubernimm sie
            \emph{negiert} (z.\,B.\ $\neg X$).
      \end{itemize}
\item Verbinde alle Minterme mit $\vee$. Fertig ist die DNF.
\end{enumerate}
\end{merkbox}

\begin{merkbox}[title={Strategie: Wahrheitstabelle $\to$ CNF}]
\textbf{CNF-Rezept:}
\begin{enumerate}[nosep]
\item Gehe jede Zeile durch, bei der die Formel den Wert $0$ hat
      ("`Nuller-Zeilen"').
\item F\"ur jede solche Zeile bildest du einen Maxterm (eine
      ODER-Verkn\"upfung), aber \textbf{umgekehrt} im Vorzeichen:
      \begin{itemize}[nosep]
      \item Ist eine Variable in der Zeile~$1$, \"ubernimm sie
            \emph{negiert} (z.\,B.\ $\neg X$).
      \item Ist eine Variable in der Zeile~$0$, \"ubernimm sie
            \emph{positiv} (z.\,B.\ $X$).
      \end{itemize}
\item Verbinde alle Maxterme mit $\wedge$. Fertig ist die CNF.
\end{enumerate}
\end{merkbox}

\begin{merkbox}
\textbf{Kurzform zum Merken:}
\begin{center}
"`\textbf{DNF aus den Einsen, CNF aus den Nullen} -- bei der CNF
dreht man zus\"atzlich die Literale um."'
\end{center}
\end{merkbox}

\begin{analogie}[title={Analogie: Einsen zusammenf\"ugen, Nullen ausschlie{\ss}en}]
\textbf{DNF-Idee:} Du hast eine Liste aller Zeilen, in denen es
"`klappt"'. Du baust f\"ur jede Zeile eine Blaupause: "`\emph{Genau
diese} Belegung der Variablen macht mich wahr."' Die DNF sagt: "`Sei
eine dieser Blaupausen wahr, dann bin ich wahr."' Deshalb ODER.

\medskip
\textbf{CNF-Idee:} Statt die "`Einsen"' aufzuz\"ahlen, schlie{\ss}t du
die "`Nullen"' aus. F\"ur jede Null-Zeile baust du eine Klausel, die
\emph{nur in genau dieser einen Zeile} falsch ist -- und in allen
anderen wahr. Da die Klausel in genau einer Zeile falsch wird,
"`schlie{\ss}t"' sie diese Zeile aus. Macht man das f\"ur alle
Null-Zeilen und verkn\"upft alle Klauseln mit UND, bleibt nur in den
Einser-Zeilen alles wahr. Deshalb funktioniert das Umdrehen der
Literale.
\end{analogie}

\textbf{Warum dreht man die Literale bei CNF um?} Das ist der Punkt,
der viele verwirrt. Schauen wir es uns genau an. Eine Klausel
$(L_1 \vee L_2 \vee L_3)$ ist \emph{nur dann} falsch, wenn
\emph{alle} drei Literale falsch sind. Wir wollen eine Klausel bauen,
die genau in \emph{einer bestimmten} Null-Zeile falsch wird. Also
m\"ussen wir die Literale so w\"ahlen, dass genau f\"ur diese eine
Zeile jedes einzelne Literal falsch wird.

\textbf{Beispiel.} Angenommen die Formel ist in Zeile $(X,Y,Z) = (1,0,1)$ falsch.
Welche Literale werden in dieser Zeile falsch?
\begin{itemize}[nosep]
\item $X$ ist in der Zeile $1$. Also ist $X$ wahr, $\neg X$ ist falsch.
      Wir w\"ahlen $\neg X$.
\item $Y$ ist in der Zeile $0$. Also ist $Y$ falsch, $\neg Y$ ist wahr.
      Wir w\"ahlen $Y$.
\item $Z$ ist in der Zeile $1$. Also w\"ahlen wir $\neg Z$.
\end{itemize}
Die Klausel $(\neg X \vee Y \vee \neg Z)$ wird in der Zeile $(1,0,1)$
zu $(0 \vee 0 \vee 0) = 0$ -- wie gew\"unscht. In jeder anderen Zeile
wird mindestens ein Literal wahr, die Klausel ist dort also wahr.
Jetzt merkt man die Regel: F\"ur eine Variable, die in der
betrachteten Zeile den Wert $1$ hat, schreibt man im Maxterm
$\neg$Variable. F\"ur eine Variable mit Wert $0$ in der Zeile schreibt
man Variable ohne Negation. \textbf{Umgekehrt zu den Einsen bei der
DNF} -- genau das ist der "`Trick"' mit dem Umdrehen der Literale.

\subsection*{0.11 Funktionale Vollst\"andigkeit}

\begin{merkbox}
Eine Menge $M$ von Operatoren (Junktoren) hei{\ss}t \textbf{funktional
vollst\"andig}, wenn man \emph{jede} boolesche Funktion durch eine
Formel ausdr\"ucken kann, die nur Operatoren aus $M$ verwendet. Kurz:
"`Mit diesen Werkzeugen kann man alles bauen."'
\end{merkbox}

\begin{analogie}
Stell dir vor, du darfst einen Werkzeugkasten zum Bau eines Hauses
mitnehmen. Ein Werkzeugkasten ist "`vollst\"andig"', wenn du damit
alles reparieren oder bauen kannst. Ein Werkzeugkasten mit nur einem
Schraubenzieher ist \emph{nicht} vollst\"andig -- du kommst nicht an
jede Ecke. Ein Werkzeugkasten mit Hammer, Zange, Schraubenzieher,
S\"age, Bohrer ist ganz sicher vollst\"andig. Die Frage in Aufgabe~2
ist: Reicht der kleine Werkzeugkasten $\{\vee, \neg\}$ aus, um alle
booleschen Funktionen zu "`bauen"'?
\end{analogie}

In Aufgabe~2 zeigen wir genau das: $\{\vee, \neg\}$ ist funktional
vollst\"andig. Das hei{\ss}t, jede Formel, die andere Operatoren
benutzt ($\wedge$, $\to$, $\leftrightarrow$), kann man so umformen,
dass nur noch $\vee$ und $\neg$ vorkommen.

\subsection*{0.12 Strukturelle Induktion}

Jetzt kommt die Beweistechnik aus Aufgabe~2. Sie hei{\ss}t
\textbf{strukturelle Induktion} und ist eigentlich ganz harmlos, wenn
man sie einmal verstanden hat.

\begin{merkbox}
\textbf{Strukturelle Induktion} ist eine Beweistechnik, mit der man
Aussagen \"uber \emph{alle} Formeln zeigen kann. Sie hat zwei
Schritte:
\begin{enumerate}[nosep]
\item \textbf{Induktionsanfang:} Zeige, dass die Aussage f\"ur die
      einfachsten Formeln gilt (die "`Bl\"atter"' des Syntaxbaums --
      meist atomare Variablen).
\item \textbf{Induktionsschritt:} Zeige, dass aus der G\"ultigkeit
      der Aussage f\"ur \emph{alle echten Teilformeln} einer Formel
      auch die G\"ultigkeit f\"ur die Gesamtformel folgt.
\end{enumerate}
\end{merkbox}

\begin{analogie}[title={Analogie: Dominosteine}]
Stell dir eine lange Reihe von Dominosteinen vor.
\begin{enumerate}[nosep]
\item \textbf{Schritt 1 -- Anfang:} Du zeigst, dass der erste Stein
      f\"allt. (Aussage gilt f\"ur die einfachsten Formeln.)
\item \textbf{Schritt 2 -- Sto{\ss}en:} Du zeigst, dass \emph{jeder}
      fallende Stein den n\"achsten umst\"o{\ss}t. (Aus G\"ultigkeit
      f\"ur Teilformeln folgt G\"ultigkeit f\"ur die Gesamtformel.)
\end{enumerate}
Folge: Alle Steine fallen -- auch ohne dass wir jeden einzelnen
anschauen m\"ussen. Die strukturelle Induktion gilt f\"ur \emph{alle}
Formeln, weil jede Formel irgendwann aus endlich vielen solcher
Schritte aufgebaut wurde.
\end{analogie}

\textbf{Warum "`strukturell"'?} Weil wir der Reihe nach alle
m\"oglichen \emph{Strukturen} (Formen) einer Formel durchgehen: Ist
die Formel eine atomare Variable? Eine Negation? Eine Konjunktion?
Eine Disjunktion? Eine Implikation? Eine Biimplikation? F\"ur jede
dieser Strukturen zeigen wir einen eigenen Fall.

\bigskip
Damit haben wir alle Werkzeuge zusammen. Jetzt k\"onnen wir uns an
die drei Aufgaben machen.

\newpage

% ============================================================
% AUFGABE 1
% ============================================================

\section*{Aufgabe 1: CNF und DNF f\"ur $\neg X \leftrightarrow Z \to \neg Y \wedge Z$}

\textbf{Aufgabenstellung in einfachen Worten.} Wir bekommen eine
Formel ohne Klammern. Daraus sollen wir zwei \"aquivalente Formeln
bauen: eine in Konjunktiver Normalform (CNF) und eine in Disjunktiver
Normalform (DNF). Der Weg daf\"ur: Zuerst verstehen, wie die Formel
geklammert wird; dann eine Wahrheitstabelle aufstellen; und am Ende
die CNF/DNF daraus ablesen.

\subsection*{Schritt 1: Die Formel klammern}

Die Formel ist
\[
F \;\equiv\; \neg X \leftrightarrow Z \to \neg Y \wedge Z.
\]
Ohne Klammern k\"onnte man sie auf viele Arten lesen. Die
Operatorpr\"azedenz aus Teil~0 sagt uns, was Vorrang hat:
$\neg > \wedge > \vee > \to > \leftrightarrow$.

Wir gehen von \emph{innen nach au{\ss}en}:

\begin{enumerate}
\item \textbf{Negationen} binden am st\"arksten. Also werden
      $\neg X$ und $\neg Y$ zu eigenen kleinen Bausteinen:
\[
\underline{\neg X} \;\leftrightarrow\; Z \to \underline{\neg Y} \wedge Z.
\]

\item \textbf{Konjunktion} ($\wedge$) bindet st\"arker als
      $\to$ und $\leftrightarrow$. Also wird $\neg Y \wedge Z$ zu
      einem Baustein:
\[
\neg X \;\leftrightarrow\; Z \to \underline{(\neg Y \wedge Z)}.
\]

\item \textbf{Implikation} ($\to$) bindet st\"arker als
      $\leftrightarrow$. Also wird die rechte Seite der
      Biimplikation zu einem Baustein:
\[
\neg X \;\leftrightarrow\; \underline{\bigl(Z \to (\neg Y \wedge Z)\bigr)}.
\]

\item Zuletzt bleibt die \textbf{Biimplikation} als \"au{\ss}erster
      Operator:
\[
F \;\equiv\; \neg X \;\leftrightarrow\; \bigl(\,Z \to (\neg Y \wedge Z)\,\bigr).
\]
\end{enumerate}

\begin{merkbox}
Die endg\"ultige Klammerung lautet:
\[
F \;\equiv\; \neg X \;\leftrightarrow\; \bigl(\,Z \to (\neg Y \wedge Z)\,\bigr).
\]
Das ist die \textbf{einzige} Lesart, die mit der Pr\"azedenzregel
vertr\"aglich ist. Alle anderen Lesarten w\"aren schlicht falsch.
\end{merkbox}

\textbf{Der Syntaxbaum anschaulich.} Man kann sich die Formel als
einen kleinen Baum vorstellen: Die Wurzel ist der oberste Operator,
darunter h\"angen die Operanden. F\"ur unsere Formel:

\begin{center}
\begin{tabular}{c}
$\leftrightarrow$ (Wurzel) \\
$\diagup \quad \diagdown$ \\
$\neg X$ \quad $\to$ \\
\quad\quad\quad $\diagup \quad \diagdown$ \\
\quad\quad\quad $Z$ \quad $\wedge$ \\
\quad\quad\quad\quad\quad\quad $\diagup \quad \diagdown$ \\
\quad\quad\quad\quad\quad\quad $\neg Y$ \quad $Z$ \\
\end{tabular}
\end{center}

Lies von unten nach oben: Ganz unten die Bl\"atter ($\neg Y$, $Z$),
dar\"uber $\wedge$, dar\"uber $\to$, dar\"uber $\leftrightarrow$, die
Wurzel. Genau diese Baumstruktur entspricht der Klammerung von oben.

\subsection*{Schritt 2: Wahrheitstabelle aufstellen}

Wir haben drei Variablen ($X$, $Y$, $Z$). Also brauchen wir
$2^3 = 8$~Zeilen. Um Fehler zu vermeiden, berechnen wir
\textbf{Spalte f\"ur Spalte} von innen nach au{\ss}en -- genau in
derselben Reihenfolge, in der wir die Formel geklammert haben.

\begin{enumerate}[nosep]
\item Erst $\neg X$ und $\neg Y$ (das sind direkt die Negationen).
\item Dann $\neg Y \wedge Z$ (Konjunktion).
\item Dann $Z \to (\neg Y \wedge Z)$ (Implikation).
\item Zuletzt $\neg X \leftrightarrow (Z \to (\neg Y \wedge Z))$
      (Biimplikation). Das ist die ganze Formel $F$.
\end{enumerate}

Hier die vollst\"andige Tabelle:

\begin{center}
\renewcommand{\arraystretch}{1.2}
\begin{tabular}{ccc|c|c|c|c|c}
\toprule
$X$ & $Y$ & $Z$ & $\neg X$ & $\neg Y$ & $\neg Y \wedge Z$ &
$Z \to (\neg Y \wedge Z)$ & $F$ \\
\midrule
0 & 0 & 0 & 1 & 1 & 0 & 1 & \textbf{1} \\
0 & 0 & 1 & 1 & 1 & 1 & 1 & \textbf{1} \\
0 & 1 & 0 & 1 & 0 & 0 & 1 & \textbf{1} \\
0 & 1 & 1 & 1 & 0 & 0 & 0 & 0 \\
1 & 0 & 0 & 0 & 1 & 0 & 1 & 0 \\
1 & 0 & 1 & 0 & 1 & 1 & 1 & 0 \\
1 & 1 & 0 & 0 & 0 & 0 & 1 & 0 \\
1 & 1 & 1 & 0 & 0 & 0 & 0 & \textbf{1} \\
\bottomrule
\end{tabular}
\end{center}

\textbf{Durchsprache einzelner Zeilen.} Wir gehen ein paar Zeilen
ausf\"uhrlich durch, damit wirklich kein Zweifel bleibt.

\medskip
\textbf{Zeile 1: $(X,Y,Z) = (0,0,0)$.}
\begin{itemize}[nosep]
\item $\neg X$: $X = 0$, also $\neg X = 1$.
\item $\neg Y$: $Y = 0$, also $\neg Y = 1$.
\item $\neg Y \wedge Z$: $1 \wedge 0 = 0$ (weil $Z=0$).
\item $Z \to (\neg Y \wedge Z)$: $0 \to 0 = 1$. Die Implikation ist
      immer wahr, wenn die Voraussetzung ($Z$) falsch ist -- erinnere
      dich an die Analogie mit dem Regenversprechen: Wenn es nicht
      regnet, kann niemand dein Versprechen brechen.
\item $F \equiv \neg X \leftrightarrow \dots$: $1 \leftrightarrow 1 = 1$.
      Beide Seiten haben denselben Wert, die Biimplikation ist wahr.
\end{itemize}

\textbf{Zeile 2: $(X,Y,Z) = (0,0,1)$.}
\begin{itemize}[nosep]
\item $\neg X = 1$.
\item $\neg Y = 1$.
\item $\neg Y \wedge Z$: $1 \wedge 1 = 1$.
\item $Z \to (\neg Y \wedge Z)$: $1 \to 1 = 1$. Hier gilt: Regenwarnung
      (Voraussetzung $Z = 1$) und Schirm dabei (Folge $= 1$). Alles
      gut.
\item $F$: $1 \leftrightarrow 1 = 1$.
\end{itemize}

\textbf{Zeile 4: $(X,Y,Z) = (0,1,1)$.} Hier wird es erstmals
interessant.
\begin{itemize}[nosep]
\item $\neg X = 1$.
\item $\neg Y = 0$ (weil $Y = 1$).
\item $\neg Y \wedge Z$: $0 \wedge 1 = 0$.
\item $Z \to (\neg Y \wedge Z)$: $1 \to 0 = 0$. \textbf{Das ist der
      eine Fall, in dem die Implikation falsch ist:} Voraussetzung
      wahr, Folge falsch. Regenwarnung, aber kein Schirm.
\item $F$: $1 \leftrightarrow 0 = 0$. Die beiden Seiten haben
      unterschiedliche Werte, die Biimplikation ist falsch.
\end{itemize}

\textbf{Zeile 8: $(X,Y,Z) = (1,1,1)$.}
\begin{itemize}[nosep]
\item $\neg X = 0$.
\item $\neg Y = 0$.
\item $\neg Y \wedge Z$: $0 \wedge 1 = 0$.
\item $Z \to (\neg Y \wedge Z)$: $1 \to 0 = 0$. Wieder der falsche
      Fall der Implikation.
\item $F$: $0 \leftrightarrow 0 = 1$. \textbf{Beide Seiten sind
      falsch} -- und das macht die Biimplikation \emph{wahr}. Sie ist
      n\"amlich nicht nur bei "`beide wahr"', sondern auch bei
      "`beide falsch"' wahr.
\end{itemize}

Die \"ubrigen Zeilen funktionieren analog. Man sollte einmal mit dem
Bleistift alle 8~Zeilen durchgehen, das sichert das Verst\"andnis.

\begin{merkbox}
\textbf{Wahre Zeilen} (hier steht $F = 1$):
\[
(X,Y,Z) \in \{(0,0,0),\,(0,0,1),\,(0,1,0),\,(1,1,1)\}.
\]
\textbf{Falsche Zeilen} (hier steht $F = 0$):
\[
(X,Y,Z) \in \{(0,1,1),\,(1,0,0),\,(1,0,1),\,(1,1,0)\}.
\]
Aus den wahren Zeilen bauen wir die DNF, aus den falschen die CNF.
\end{merkbox}

\subsection*{Schritt 3: DNF ablesen}

Wir gehen jetzt die \textbf{4 wahren Zeilen} durch und bauen f\"ur
jede einen Minterm. Rezept: Wo die Variable in der Zeile $1$ ist,
schreiben wir sie positiv; wo sie $0$ ist, schreiben wir sie negiert.

\textbf{Zeile 1: $(X,Y,Z) = (0,0,0)$.}
\begin{itemize}[nosep]
\item $X = 0 \Rightarrow$ $\neg X$
\item $Y = 0 \Rightarrow$ $\neg Y$
\item $Z = 0 \Rightarrow$ $\neg Z$
\end{itemize}
Minterm: $\neg X \wedge \neg Y \wedge \neg Z$. Dieser Minterm ist
\emph{nur} in der Zeile $(0,0,0)$ wahr (da sind alle drei negierten
Variablen wahr) und in allen anderen Zeilen falsch.

\textbf{Zeile 2: $(X,Y,Z) = (0,0,1)$.}
\begin{itemize}[nosep]
\item $X = 0 \Rightarrow$ $\neg X$
\item $Y = 0 \Rightarrow$ $\neg Y$
\item $Z = 1 \Rightarrow$ $Z$
\end{itemize}
Minterm: $\neg X \wedge \neg Y \wedge Z$.

\textbf{Zeile 3: $(X,Y,Z) = (0,1,0)$.}
Minterm: $\neg X \wedge Y \wedge \neg Z$.

\textbf{Zeile 8: $(X,Y,Z) = (1,1,1)$.}
Minterm: $X \wedge Y \wedge Z$.

Alle 4 Minterme werden mit ODER verkn\"upft. Das ergibt die DNF:

\begin{merkbox}[title={Ergebnis: DNF}]
\[
\text{DNF}(F) \;\equiv\; \begin{aligned}[t]
  &(\neg X \wedge \neg Y \wedge \neg Z) \;\vee\; \\
  &(\neg X \wedge \neg Y \wedge Z) \;\vee\; \\
  &(\neg X \wedge Y \wedge \neg Z) \;\vee\; \\
  &(X \wedge Y \wedge Z).
\end{aligned}
\]
\end{merkbox}

\subsection*{Schritt 4: CNF ablesen}

Wir gehen jetzt die \textbf{4 falschen Zeilen} durch und bauen f\"ur
jede einen Maxterm. Rezept: Wo die Variable in der Zeile $1$ ist,
schreiben wir sie \emph{negiert}; wo sie $0$ ist, schreiben wir sie
\emph{positiv}. (Umgekehrt zu DNF!)

\textbf{Zeile 4: $(X,Y,Z) = (0,1,1)$.}
\begin{itemize}[nosep]
\item $X = 0 \Rightarrow$ positiv: $X$
\item $Y = 1 \Rightarrow$ negiert: $\neg Y$
\item $Z = 1 \Rightarrow$ negiert: $\neg Z$
\end{itemize}
Maxterm: $(X \vee \neg Y \vee \neg Z)$. Diese Klausel wird \emph{nur}
in der Zeile $(0,1,1)$ falsch: dort ist $X = 0$, $\neg Y = 0$,
$\neg Z = 0$, also $0 \vee 0 \vee 0 = 0$. In allen anderen Zeilen ist
mindestens eines der drei Literale wahr.

\textbf{Zeile 5: $(X,Y,Z) = (1,0,0)$.}
\begin{itemize}[nosep]
\item $X = 1 \Rightarrow$ $\neg X$
\item $Y = 0 \Rightarrow$ $Y$
\item $Z = 0 \Rightarrow$ $Z$
\end{itemize}
Maxterm: $(\neg X \vee Y \vee Z)$.

\textbf{Zeile 6: $(X,Y,Z) = (1,0,1)$.}
\begin{itemize}[nosep]
\item $X = 1 \Rightarrow$ $\neg X$
\item $Y = 0 \Rightarrow$ $Y$
\item $Z = 1 \Rightarrow$ $\neg Z$
\end{itemize}
Maxterm: $(\neg X \vee Y \vee \neg Z)$.

\textbf{Zeile 7: $(X,Y,Z) = (1,1,0)$.}
\begin{itemize}[nosep]
\item $X = 1 \Rightarrow$ $\neg X$
\item $Y = 1 \Rightarrow$ $\neg Y$
\item $Z = 0 \Rightarrow$ $Z$
\end{itemize}
Maxterm: $(\neg X \vee \neg Y \vee Z)$.

Alle 4 Maxterme werden mit UND verkn\"upft. Das ergibt die CNF:

\begin{merkbox}[title={Ergebnis: CNF}]
\[
\text{CNF}(F) \;\equiv\; \begin{aligned}[t]
  &(X \vee \neg Y \vee \neg Z) \;\wedge\; \\
  &(\neg X \vee Y \vee Z) \;\wedge\; \\
  &(\neg X \vee Y \vee \neg Z) \;\wedge\; \\
  &(\neg X \vee \neg Y \vee Z).
\end{aligned}
\]
\end{merkbox}

\subsection*{Schritt 5: Warum das Umdrehen bei CNF funktioniert}

Das ist der Punkt, der oft fehlt. Wir haben es schon in Teil~0
erkl\"art, aber hier nochmal an einem konkreten Beispiel.

Nimm die erste Klausel der CNF: $(X \vee \neg Y \vee \neg Z)$. In
welchen Zeilen ist sie falsch?
\begin{itemize}[nosep]
\item Damit die Disjunktion $(X \vee \neg Y \vee \neg Z)$ falsch wird,
      m\"ussen alle drei Literale falsch sein.
\item $X$ falsch bedeutet $X = 0$.
\item $\neg Y$ falsch bedeutet $Y = 1$.
\item $\neg Z$ falsch bedeutet $Z = 1$.
\end{itemize}
Die Klausel ist also \emph{nur} in der Zeile $(X,Y,Z) = (0,1,1)$
falsch. In allen anderen sieben Zeilen ist sie wahr.

Das ist \emph{exakt} die Zeile, f\"ur die wir sie gebaut haben. Sie
"`schlie{\ss}t"' also nur diese eine Zeile aus. Die zweite Klausel
schlie{\ss}t nur $(1,0,0)$ aus, die dritte nur $(1,0,1)$, die vierte
nur $(1,1,0)$. Die Gesamte CNF ist die Konjunktion dieser vier
Klauseln. Sie ist \emph{nur} dann wahr, wenn \emph{keine} dieser vier
Zeilen getroffen wird -- also wenn wir in einer der vier
\emph{wahren} Zeilen sind. Genau wie die Formel $F$ selbst. Damit ist
die CNF zu $F$ \"aquivalent. Genau was wir wollten!

\subsection*{Zusammenfassung Aufgabe~1}

\begin{merkbox}[title={Ergebnis Aufgabe~1}]
F\"ur $F \;\equiv\; \neg X \leftrightarrow (Z \to (\neg Y \wedge Z))$
gilt:

\medskip
\textbf{DNF:}
\[
(\neg X \wedge \neg Y \wedge \neg Z)
 \vee (\neg X \wedge \neg Y \wedge Z)
 \vee (\neg X \wedge Y \wedge \neg Z)
 \vee (X \wedge Y \wedge Z).
\]

\textbf{CNF:}
\[
(X \vee \neg Y \vee \neg Z)
 \wedge (\neg X \vee Y \vee Z)
 \wedge (\neg X \vee Y \vee \neg Z)
 \wedge (\neg X \vee \neg Y \vee Z).
\]

\medskip
Sanity-Check: 4~wahre Zeilen + 4~falsche Zeilen = 8~Gesamtzeilen.
Passt.
\end{merkbox}

\newpage

% ============================================================
% AUFGABE 2
% ============================================================

\section*{Aufgabe 2: Jede Formel ist \"aquivalent zu einer nur mit $\vee$ und $\neg$}

\textbf{Aufgabenstellung in einfachen Worten.} Wir sollen zeigen:
Egal wie kompliziert eine logische Formel ist -- wie viele $\wedge$,
$\to$ oder $\leftrightarrow$ sie auch enth\"alt -- man kann sie immer
in eine \"aquivalente Formel umschreiben, die nur $\vee$ und $\neg$
benutzt. Das hei{\ss}t: $\{\vee, \neg\}$ ist ein "`vollst\"andiges
Werkzeugset"'.

\begin{analogie}
Stell dir vor, du hast in der K\"uche nur ein Messer und einen
L\"offel. Kann man damit alles essen? Vielleicht mit etwas M\"uhe,
aber ja. Genauso die Frage in dieser Aufgabe: Reichen uns die beiden
Werkzeuge $\vee$ und $\neg$, um jede logische Aussage auszudr\"ucken?
Die Antwort ist \textbf{ja}, und wir zeigen das hier sauber.
\end{analogie}

\subsection*{Ziel pr\"azise formulieren}

\textbf{Behauptung.} F\"ur jede aussagenlogische Formel $F$ existiert
eine Formel $F'$ mit folgenden Eigenschaften:
\begin{enumerate}[nosep]
\item $F'$ verwendet nur die Junktoren $\vee$ und $\neg$ (also keine
      $\wedge$, $\to$, $\leftrightarrow$).
\item $F' \equiv F$, das hei{\ss}t: $F'$ hat dieselbe Wahrheitstabelle
      wie $F$.
\end{enumerate}

\subsection*{Beweisidee: Strukturelle Induktion}

Wir gehen alle m\"oglichen Formen einer Formel durch. Am Anfang steht
die einfachste: atomare Variablen. Dann kommen die Formen mit
Operatoren. F\"ur jede Form zeigen wir, wie man die Formel mithilfe
von $\vee$ und $\neg$ ausdr\"uckt.

\begin{analogie}[title={Analogie: Dominokette}]
Denk an Dominosteine: Wir stellen eine Kette auf, und jeder fallende
Stein sto{\ss}t den n\"achsten um. Hier sind die Steine die
verschiedenen M\"oglichkeiten, wie eine Formel aufgebaut sein kann.
Wir zeigen: Der erste Stein f\"allt (Atome sind trivial), und jeder
zusammengesetzte Stein f\"allt, wenn seine Teilformeln schon fallen
k\"onnen. Das l\"auft bis ganz oben zum Wurzeloperator jeder Formel.
\end{analogie}

\subsection*{Induktionsanfang: atomare Formeln}

Sei $F$ eine \textbf{atomare Formel}, also nur eine einzelne
Aussagenvariable $p$ (oder $q$, $r$, $X$ -- egal wie sie hei{\ss}t).
Dann enth\"alt $F$ \emph{keine Junktoren} -- also sicher keine
verbotenen $\wedge$, $\to$, $\leftrightarrow$.

Wir setzen $F' := p$. Dann ist $F' \equiv F$ trivialerweise (sie sind
ja dieselbe Formel), und $F'$ verwendet nur... naja, eigentlich gar
keine Junktoren. Das ist aber erlaubt: "`nur $\vee$ und $\neg$
verwenden"' schlie{\ss}t nicht aus, dass man \emph{gar nichts}
verwendet.

\textbf{Fertig mit dem Anfang.} Der erste Dominostein ist umgefallen.

\subsection*{Induktionsschritt: zusammengesetzte Formeln}

Sei nun $F$ eine zusammengesetzte Formel (also nicht nur eine
einzelne Variable). Dann hat $F$ einen \textbf{obersten Operator}
("`Hauptoperator"'). Je nach diesem Operator hat $F$ eine
bestimmte Form.

\textbf{Induktionsvoraussetzung (IV).} F\"ur alle \emph{echten
Teilformeln} von $F$ gilt die Behauptung bereits. Das hei{\ss}t:
Zu jeder Teilformel $G$ von $F$ existiert schon eine
\"aquivalente Formel $G'$ \"uber $\{\vee, \neg\}$. Das k\"onnen wir
annehmen, weil Teilformeln strikt kleiner sind als $F$, und f\"ur
kleinere Formeln haben wir die Behauptung als Induktionsannahme.

Wir unterscheiden jetzt nach dem Hauptoperator von $F$. Es gibt f\"unf
F\"alle.

\subsubsection*{Fall 1: $F = \neg G$}

Die Formel ist eine Negation, der Hauptoperator ist also $\neg$. Nach
IV existiert $G' \equiv G$ mit $G' \in \{\vee, \neg\}$. Setze
\[
F' \;:=\; \neg G'.
\]
Dann ist $F' = \neg G' \equiv \neg G = F$. Und $F'$ benutzt nur
$\vee$ und $\neg$ (die Negation geh\"ort dazu, und $G'$ auch).
\textbf{Fertig.}

\subsubsection*{Fall 2: $F = G \vee H$}

Die Formel ist eine Disjunktion. Nach IV existieren
$G' \equiv G$ und $H' \equiv H$ mit $G', H' \in \{\vee, \neg\}$.
Setze
\[
F' \;:=\; G' \vee H'.
\]
Dann ist $F' \equiv G \vee H = F$, und $F'$ benutzt nur $\vee$ und
$\neg$. \textbf{Fertig.}

\subsubsection*{Fall 3: $F = G \wedge H$}

Jetzt wird es interessant: Die Formel benutzt ein $\wedge$, das wir
"`wegbekommen"' m\"ussen. Wir verwenden die \textbf{De-Morgan-Regel}:
\[
G \wedge H \;\equiv\; \neg(\neg G \vee \neg H).
\]

\textbf{Warum gilt das?} \"Uberlegen wir in Alltagssprache:
\begin{itemize}[nosep]
\item $G \wedge H$ ist genau dann wahr, wenn \emph{beide}, $G$ und
      $H$, wahr sind -- wenn also \emph{keiner von beiden falsch}
      ist.
\item "`Keiner von beiden falsch"' bedeutet: $\neg G$ ist falsch UND
      $\neg H$ ist falsch.
\item Wenn $\neg G$ und $\neg H$ beide falsch sind, ist auch ihre
      Disjunktion $\neg G \vee \neg H$ falsch.
\item Und wenn $\neg G \vee \neg H$ falsch ist, dann ist
      $\neg(\neg G \vee \neg H)$ wahr.
\end{itemize}
Ergebnis: $G \wedge H$ wahr $\iff$ $\neg(\neg G \vee \neg H)$ wahr.
Also sind die beiden Formeln \"aquivalent. \checkmark

Nach IV existieren $G' \equiv G$ und $H' \equiv H$ \"uber
$\{\vee, \neg\}$. Setze
\[
F' \;:=\; \neg(\neg G' \vee \neg H').
\]
Dann ist $F' \equiv \neg(\neg G \vee \neg H) \equiv G \wedge H = F$,
und $F'$ benutzt nur $\vee$ und $\neg$. \textbf{Fertig.}

\subsubsection*{Fall 4: $F = G \to H$}

Jetzt haben wir eine Implikation, die auch "`weg"' muss. Die Regel
ist:
\[
G \to H \;\equiv\; \neg G \vee H.
\]

\textbf{Warum gilt das?} Die Implikation $G \to H$ ist (wie wir aus
Teil~0 wissen) genau dann falsch, wenn $G$ wahr und $H$ falsch.
Schauen wir, wann $\neg G \vee H$ falsch ist:
\begin{itemize}[nosep]
\item $\neg G \vee H$ ist falsch, wenn \emph{beide} Disjunkte falsch
      sind.
\item $\neg G$ falsch bedeutet: $G$ wahr.
\item $H$ falsch bedeutet: $H$ falsch.
\end{itemize}
Also ist $\neg G \vee H$ falsch genau dann, wenn $G$ wahr und $H$
falsch ist -- dieselbe Bedingung wie f\"ur $G \to H$. In allen anderen
drei F\"allen sind beide Formeln wahr. Die Wahrheitstabellen stimmen
\"uberein. \checkmark

\begin{center}
\renewcommand{\arraystretch}{1.2}
\begin{tabular}{cc|c|c|c}
\toprule
$G$ & $H$ & $\neg G$ & $\neg G \vee H$ & $G \to H$ \\
\midrule
0 & 0 & 1 & 1 & 1 \\
0 & 1 & 1 & 1 & 1 \\
1 & 0 & 0 & 0 & 0 \\
1 & 1 & 0 & 1 & 1 \\
\bottomrule
\end{tabular}
\end{center}

Die beiden rechten Spalten sind identisch -- damit ist die
\"Aquivalenz gezeigt.

Nach IV existieren $G' \equiv G$ und $H' \equiv H$. Setze
\[
F' \;:=\; \neg G' \vee H'.
\]
\textbf{Fertig.}

\subsubsection*{Fall 5: $F = G \leftrightarrow H$}

Die Biimplikation ist der komplizierteste Fall, aber wir zerlegen
sie Schritt f\"ur Schritt:
\[
G \leftrightarrow H \;\equiv\; (G \to H) \wedge (H \to G).
\]

\textbf{Warum gilt das?} Die Biimplikation sagt "`beide haben
denselben Wert"', und das ist genau dasselbe wie "`$G$ impliziert $H$
\emph{und} $H$ impliziert $G$"'. Anders gesagt: Wenn beide Implikationen
gelten, kann man weder in $G$ noch in $H$ sein, ohne auch im anderen
zu sein. Beide m\"ussen also gleich wahr/falsch sein.

Nun ersetzen wir die Implikationen nach Fall~4:
\[
(G \to H) \wedge (H \to G) \;\equiv\; (\neg G \vee H) \wedge (\neg H \vee G).
\]
Wir sind fast da, aber es ist noch ein $\wedge$ \"ubrig. Das ersetzen
wir nach Fall~3 (De Morgan):
\[
(\neg G \vee H) \wedge (\neg H \vee G) \;\equiv\;
\neg\bigl(\neg(\neg G \vee H) \vee \neg(\neg H \vee G)\bigr).
\]
Jetzt sind keine $\wedge$, $\to$ oder $\leftrightarrow$ mehr
enthalten. Nach IV ersetzen wir $G \to G'$ und $H \to H'$:
\[
F' \;:=\; \neg\bigl(\neg(\neg G' \vee H') \vee \neg(\neg H' \vee G')\bigr).
\]
\textbf{Fertig.}

\subsection*{Abschluss des Beweises}

Wir haben f\"ur jeden m\"oglichen Hauptoperator gezeigt, wie man die
Formel in eine \"aquivalente $\vee/\neg$-Formel \"uberf\"uhrt, unter
Verwendung der IV f\"ur die Teilformeln. Da jede Formel einen
\emph{endlichen} Syntaxbaum hat, k\"onnen wir den Prozess immer
konsequent von den Bl\"attern zur Wurzel hochklettern: Zuerst l\"osen
wir die atomaren Formeln (Induktionsanfang), dann die Operatoren
auf der Ebene dar\"uber, dann die n\"achste Ebene, und so weiter, bis
wir oben bei der Wurzel angekommen sind. Die Umwandlung terminiert
wegen der Endlichkeit des Syntaxbaums, und das Ergebnis ist eine
\"aquivalente Formel, die nur $\vee$ und $\neg$ verwendet.

\textbf{Damit ist gezeigt: Jede Formel ist \"aquivalent zu einer
Formel \"uber $\{\vee, \neg\}$.} $\qquad\square$

\begin{merkbox}[title={Ergebnis Aufgabe~2: $\{\vee, \neg\}$ ist funktional vollst\"andig}]
Die Ersetzungsregeln aus dem Beweis sind:
\begin{align*}
G \wedge H &\;\equiv\; \neg(\neg G \vee \neg H) \\
G \to H &\;\equiv\; \neg G \vee H \\
G \leftrightarrow H &\;\equiv\;
   \neg\bigl(\neg(\neg G \vee H) \vee \neg(\neg H \vee G)\bigr)
\end{align*}
Mit diesen drei Regeln und der Doppelnegation kann man jede Formel
Schritt f\"ur Schritt in $\vee/\neg$-Form bringen.
\end{merkbox}

\begin{merkbox}[title={Bonus: weitere funktional vollst\"andige Mengen}]
$\{\vee, \neg\}$ ist nicht die einzige funktional vollst\"andige
Menge. Auch diese reichen:
\begin{itemize}[nosep]
\item $\{\wedge, \neg\}$ -- genau dieselbe Idee, nur mit vertauschten
      Rollen.
\item $\{\to, \neg\}$ -- Implikation plus Negation reichen auch.
\item $\{\to, \bot\}$ -- Implikation plus "`falsch"' als Konstante.
\item $\{\text{NAND}\}$ -- der Sheffer-Strich ist allein vollst\"andig!
\item $\{\text{NOR}\}$ -- das Peirce-Symbol ist auch allein
      vollst\"andig.
\end{itemize}
NAND und NOR sind in der Digitaltechnik besonders wichtig, weil man
mit nur einem Gattertyp jedes beliebige Schaltwerk bauen kann.
\end{merkbox}

\newpage

% ============================================================
% AUFGABE 3
% ============================================================

\section*{Aufgabe 3: CNF und DNF aus einer Wahrheitstabelle}

\textbf{Aufgabenstellung in einfachen Worten.} Wir bekommen eine
Wahrheitstabelle mit 16 Zeilen (also f\"ur vier Variablen $A, B, C, D$).
Daraus sollen wir die CNF und die DNF der zugeh\"origen Formel
schreiben. Der Weg ist genau derselbe wie in Aufgabe~1, nur mit
einer gr\"o{\ss}eren Tabelle.

\subsection*{Strategie wiederholen}

\begin{merkbox}
\textbf{DNF:} Disjunktion (ODER) aller Minterme f\"ur die Zeilen, in
denen die Formel wahr ist (Einser-Zeilen). Literale wie in der Zeile
stehen: $1 \to $~Variable, $0 \to$~negierte Variable.

\medskip
\textbf{CNF:} Konjunktion (UND) aller Maxterme f\"ur die Zeilen, in
denen die Formel falsch ist (Nuller-Zeilen). Literale \textbf{umgedreht}:
$1 \to$~negierte Variable, $0 \to$~Variable.
\end{merkbox}

\subsection*{Schritt 1: Zeilen z\"ahlen und markieren}

Wir schreiben die Wahrheitstabelle nochmal sauber auf und markieren
neben jeder Zeile, ob sie f\"ur CNF oder DNF verwendet wird.

\begin{center}
\renewcommand{\arraystretch}{1.1}
\begin{tabular}{r|cccc|c|l}
\toprule
\# & $A$ & $B$ & $C$ & $D$ & $f$ & verwendet f\"ur \\
\midrule
 1 & 0 & 0 & 0 & 0 & 0 & CNF-Maxterm \\
 2 & 0 & 0 & 0 & 1 & \textbf{1} & DNF-Minterm \\
 3 & 0 & 0 & 1 & 0 & \textbf{1} & DNF-Minterm \\
 4 & 0 & 0 & 1 & 1 & 0 & CNF-Maxterm \\
 5 & 0 & 1 & 0 & 0 & \textbf{1} & DNF-Minterm \\
 6 & 0 & 1 & 0 & 1 & \textbf{1} & DNF-Minterm \\
 7 & 0 & 1 & 1 & 0 & 0 & CNF-Maxterm \\
 8 & 0 & 1 & 1 & 1 & \textbf{1} & DNF-Minterm \\
 9 & 1 & 0 & 0 & 0 & 0 & CNF-Maxterm \\
10 & 1 & 0 & 0 & 1 & 0 & CNF-Maxterm \\
11 & 1 & 0 & 1 & 0 & \textbf{1} & DNF-Minterm \\
12 & 1 & 0 & 1 & 1 & 0 & CNF-Maxterm \\
13 & 1 & 1 & 0 & 0 & \textbf{1} & DNF-Minterm \\
14 & 1 & 1 & 0 & 1 & 0 & CNF-Maxterm \\
15 & 1 & 1 & 1 & 0 & \textbf{1} & DNF-Minterm \\
16 & 1 & 1 & 1 & 1 & \textbf{1} & DNF-Minterm \\
\bottomrule
\end{tabular}
\end{center}

\textbf{Z\"ahlen wir nach:}
\begin{itemize}[nosep]
\item Einser-Zeilen (Formel wahr): 2, 3, 5, 6, 8, 11, 13, 15, 16
      $\Rightarrow$ \textbf{9 Zeilen}.
\item Nuller-Zeilen (Formel falsch): 1, 4, 7, 9, 10, 12, 14
      $\Rightarrow$ \textbf{7 Zeilen}.
\end{itemize}
Also: $9 + 7 = 16$, passt zur Gesamtanzahl der Zeilen. Die DNF wird
9~Minterme haben, die CNF wird 7~Maxterme haben.

\subsection*{Schritt 2: DNF aufstellen}

Wir gehen die 9~Einser-Zeilen durch. F\"ur jede bauen wir einen
Minterm nach dem Rezept "`Variable wie in der Zeile"'. Wir machen die
ersten drei richtig ausf\"uhrlich und fassen den Rest knapper.

\textbf{Zeile 2: $(A,B,C,D) = (0,0,0,1)$.}
\begin{itemize}[nosep]
\item $A = 0 \Rightarrow \neg A$
\item $B = 0 \Rightarrow \neg B$
\item $C = 0 \Rightarrow \neg C$
\item $D = 1 \Rightarrow D$
\end{itemize}
Minterm: $\neg A \wedge \neg B \wedge \neg C \wedge D$.

\textbf{Zeile 3: $(A,B,C,D) = (0,0,1,0)$.}
\begin{itemize}[nosep]
\item $A = 0 \Rightarrow \neg A$
\item $B = 0 \Rightarrow \neg B$
\item $C = 1 \Rightarrow C$
\item $D = 0 \Rightarrow \neg D$
\end{itemize}
Minterm: $\neg A \wedge \neg B \wedge C \wedge \neg D$.

\textbf{Zeile 5: $(A,B,C,D) = (0,1,0,0)$.}
\begin{itemize}[nosep]
\item $A = 0 \Rightarrow \neg A$
\item $B = 1 \Rightarrow B$
\item $C = 0 \Rightarrow \neg C$
\item $D = 0 \Rightarrow \neg D$
\end{itemize}
Minterm: $\neg A \wedge B \wedge \neg C \wedge \neg D$.

\textbf{Zeilen 6, 8, 11, 13, 15, 16.} Wir gehen sie jetzt schneller
durch:
\begin{itemize}[nosep]
\item Zeile 6: $(0,1,0,1)$ $\Rightarrow$ $\neg A \wedge B \wedge \neg C \wedge D$
\item Zeile 8: $(0,1,1,1)$ $\Rightarrow$ $\neg A \wedge B \wedge C \wedge D$
\item Zeile 11: $(1,0,1,0)$ $\Rightarrow$ $A \wedge \neg B \wedge C \wedge \neg D$
\item Zeile 13: $(1,1,0,0)$ $\Rightarrow$ $A \wedge B \wedge \neg C \wedge \neg D$
\item Zeile 15: $(1,1,1,0)$ $\Rightarrow$ $A \wedge B \wedge C \wedge \neg D$
\item Zeile 16: $(1,1,1,1)$ $\Rightarrow$ $A \wedge B \wedge C \wedge D$
\end{itemize}

Alle 9 Minterme werden mit $\vee$ verkn\"upft:

\begin{merkbox}[title={Ergebnis: DNF}]
\begin{align*}
\text{DNF}(f) \;\equiv\; &(\neg A \wedge \neg B \wedge \neg C \wedge D) \vee {} \\
& (\neg A \wedge \neg B \wedge C \wedge \neg D) \vee {} \\
& (\neg A \wedge B \wedge \neg C \wedge \neg D) \vee {} \\
& (\neg A \wedge B \wedge \neg C \wedge D) \vee {} \\
& (\neg A \wedge B \wedge C \wedge D) \vee {} \\
& (A \wedge \neg B \wedge C \wedge \neg D) \vee {} \\
& (A \wedge B \wedge \neg C \wedge \neg D) \vee {} \\
& (A \wedge B \wedge C \wedge \neg D) \vee {} \\
& (A \wedge B \wedge C \wedge D).
\end{align*}
\end{merkbox}

\subsection*{Schritt 3: CNF aufstellen}

Jetzt die 7~Nuller-Zeilen. F\"ur jede bauen wir einen Maxterm mit dem
\emph{umgekehrten} Rezept: Wo die Variable in der Zeile $1$ ist,
schreiben wir sie \textbf{negiert}; wo sie $0$ ist, schreiben wir sie
\textbf{positiv}. Auch hier die ersten drei ausf\"uhrlich.

\textbf{Zeile 1: $(A,B,C,D) = (0,0,0,0)$.}
\begin{itemize}[nosep]
\item $A = 0 \Rightarrow A$ (positiv, weil Zeilenwert 0)
\item $B = 0 \Rightarrow B$
\item $C = 0 \Rightarrow C$
\item $D = 0 \Rightarrow D$
\end{itemize}
Maxterm: $(A \vee B \vee C \vee D)$. Diese Klausel wird nur in der
Zeile $(0,0,0,0)$ falsch -- da sind alle Literale $0$. In allen
anderen Zeilen ist mindestens eine der vier Variablen $1$, also wahr.

\textbf{Zeile 4: $(A,B,C,D) = (0,0,1,1)$.}
\begin{itemize}[nosep]
\item $A = 0 \Rightarrow A$
\item $B = 0 \Rightarrow B$
\item $C = 1 \Rightarrow \neg C$
\item $D = 1 \Rightarrow \neg D$
\end{itemize}
Maxterm: $(A \vee B \vee \neg C \vee \neg D)$. Kurzer Check: In der
Zeile $(0,0,1,1)$ wird das zu $(0 \vee 0 \vee \neg 1 \vee \neg 1) =
(0 \vee 0 \vee 0 \vee 0) = 0$. Passt.

\textbf{Zeile 7: $(A,B,C,D) = (0,1,1,0)$.}
\begin{itemize}[nosep]
\item $A = 0 \Rightarrow A$
\item $B = 1 \Rightarrow \neg B$
\item $C = 1 \Rightarrow \neg C$
\item $D = 0 \Rightarrow D$
\end{itemize}
Maxterm: $(A \vee \neg B \vee \neg C \vee D)$.

\textbf{Zeilen 9, 10, 12, 14.} Schneller durch:
\begin{itemize}[nosep]
\item Zeile 9: $(1,0,0,0)$ $\Rightarrow$ $(\neg A \vee B \vee C \vee D)$
\item Zeile 10: $(1,0,0,1)$ $\Rightarrow$ $(\neg A \vee B \vee C \vee \neg D)$
\item Zeile 12: $(1,0,1,1)$ $\Rightarrow$ $(\neg A \vee B \vee \neg C \vee \neg D)$
\item Zeile 14: $(1,1,0,1)$ $\Rightarrow$ $(\neg A \vee \neg B \vee C \vee \neg D)$
\end{itemize}

Alle 7 Maxterme werden mit $\wedge$ verkn\"upft:

\begin{merkbox}[title={Ergebnis: CNF}]
\begin{align*}
\text{CNF}(f) \;\equiv\; &(A \vee B \vee C \vee D) \wedge {} \\
& (A \vee B \vee \neg C \vee \neg D) \wedge {} \\
& (A \vee \neg B \vee \neg C \vee D) \wedge {} \\
& (\neg A \vee B \vee C \vee D) \wedge {} \\
& (\neg A \vee B \vee C \vee \neg D) \wedge {} \\
& (\neg A \vee B \vee \neg C \vee \neg D) \wedge {} \\
& (\neg A \vee \neg B \vee C \vee \neg D).
\end{align*}
\end{merkbox}

\subsection*{Schritt 4: Sanity-Check}

\textbf{Anzahlen stimmen:} 9~Minterme (DNF) + 7~Maxterme (CNF) =
16~Zeilen. Alles korrekt gez\"ahlt.

\textbf{Stichprobe DNF.} Zeile 2 ist $(0,0,0,1)$ mit $f = 1$. Der
entsprechende Minterm ist $\neg A \wedge \neg B \wedge \neg C \wedge D$.
Setzen wir die Werte ein:
$(\neg 0) \wedge (\neg 0) \wedge (\neg 0) \wedge 1 =
1 \wedge 1 \wedge 1 \wedge 1 = 1$. Der Minterm ist in genau dieser
Zeile wahr, wie gew\"unscht. In allen anderen Zeilen ist mindestens
ein Literal falsch (weil $(A,B,C,D)$ einen anderen Wert hat), und
damit der ganze Minterm falsch.

\textbf{Stichprobe CNF.} Zeile 9 ist $(1,0,0,0)$ mit $f = 0$. Der
entsprechende Maxterm ist $(\neg A \vee B \vee C \vee D)$. Einsetzen:
$(\neg 1 \vee 0 \vee 0 \vee 0) = (0 \vee 0 \vee 0 \vee 0) = 0$. Der
Maxterm ist in genau dieser Zeile falsch, wie gew\"unscht. In allen
anderen Zeilen ist mindestens ein Literal wahr.

\begin{merkbox}[title={Ergebnis Aufgabe~3}]
Die CNF hat 7 Klauseln (aus den 7 Nuller-Zeilen), die DNF hat
9 Minterme (aus den 9 Einser-Zeilen). Beide sind zur urspr\"unglichen
Formel \"aquivalent. Check: $7 + 9 = 16$ Zeilen.
\end{merkbox}

\newpage

% ============================================================
% ZUSAMMENFASSUNG
% ============================================================

\section*{Zusammenfassung und goldene Regeln}

\subsection*{Kurz\"ubersicht der Ergebnisse}

\begin{center}
\renewcommand{\arraystretch}{1.3}
\begin{tabular}{l|l}
\toprule
\textbf{Aufgabe} & \textbf{Ergebnis} \\
\midrule
1 (DNF) & 4 Minterme aus 4 wahren Zeilen \\
1 (CNF) & 4 Maxterme aus 4 falschen Zeilen \\
2 & $\{\vee, \neg\}$ ist funktional vollst\"andig (Induktion) \\
3 (DNF) & 9 Minterme aus 9 wahren Zeilen \\
3 (CNF) & 7 Maxterme aus 7 falschen Zeilen \\
\bottomrule
\end{tabular}
\end{center}

\subsection*{Goldene Regeln}

\begin{merkbox}[title={Goldene Regeln zum Mitnehmen}]
\begin{enumerate}
\item \textbf{Pr\"azedenz zuerst.} Bevor du eine klammerfreie Formel
      auswertest, klammere sie mental nach $\neg > \wedge > \vee > \to
      > \leftrightarrow$ (stark $\to$ schwach). Sonst passieren
      Fl\"uchtigkeitsfehler.

\item \textbf{DNF aus den Einsen.} F\"ur jede Zeile mit $f = 1$:
      Minterm bauen. Variable mit Wert $1$ $\to$ positiv; Variable mit
      Wert $0$ $\to$ negiert. Alles mit $\vee$ verbinden.

\item \textbf{CNF aus den Nullen mit umgedrehten Literalen.} F\"ur
      jede Zeile mit $f = 0$: Maxterm bauen. Variable mit Wert $1$
      $\to$ negiert; Variable mit Wert $0$ $\to$ positiv. Alles mit
      $\wedge$ verbinden.

\item \textbf{Sanity-Check.} Anzahl Minterme (DNF) + Anzahl Maxterme
      (CNF) = Anzahl Zeilen der Wahrheitstabelle.

\item \textbf{Funktionale Vollst\"andigkeit.} Um zu zeigen, dass eine
      Junktorenmenge vollst\"andig ist, reicht es, jeden "`fremden"'
      Junktor durch eine \"aquivalente Formel aus der erlaubten Menge
      zu ersetzen. F\"ur $\{\vee, \neg\}$ sind die Ersetzungen:
      \begin{align*}
        G \wedge H &\equiv \neg(\neg G \vee \neg H) \\
        G \to H &\equiv \neg G \vee H \\
        G \leftrightarrow H &\equiv \neg(\neg(\neg G \vee H) \vee \neg(\neg H \vee G))
      \end{align*}

\item \textbf{Strukturelle Induktion.} Beweist Aussagen "`f\"ur alle
      Formeln"'. Zwei Schritte: (a) Aussage gilt f\"ur atomare Formeln
      (Induktionsanfang); (b) wenn sie f\"ur alle echten Teilformeln
      gilt, gilt sie auch f\"ur die Gesamtformel (Induktionsschritt),
      und zwar getrennt f\"ur jeden m\"oglichen Hauptoperator.
\end{enumerate}
\end{merkbox}

\begin{analogie}[title={Merke: DNF vs.\ CNF}]
\textbf{DNF = Einsen aufsammeln.} F\"ur jede Zeile, in der die Formel
wahr ist, baue ein Szenario, das genau diese Zeile beschreibt. Die
DNF sagt: "`Sei eines dieser Szenarien wahr, dann bin ich wahr."'

\medskip
\textbf{CNF = Nullen ausschlie{\ss}en.} F\"ur jede Zeile, in der die
Formel falsch ist, baue eine Ausschlussklausel, die nur f\"ur diese
eine Zeile falsch wird. Die CNF sagt: "`Sei keine dieser Nuller-Zeilen
getroffen, dann bin ich wahr."'

\medskip
Beide Wege f\"uhren zum selben Ergebnis -- sie sind nur zwei
verschiedene Perspektiven auf dieselbe Funktion.
\end{analogie}

\bigskip
\bigskip

Damit sind alle drei Aufgaben inhaltlich abgeschlossen. Die
schrittweise Erkl\"arung sollte es auch jemandem, der vorher noch nie
etwas von CNF, DNF oder struktureller Induktion geh\"ort hat,
m\"oglich gemacht haben, die Gedankeng\"ange nachzuvollziehen. Viel
Spa{\ss} beim Nachrechnen!

\end{document}
